\chapter{Literature Survey}
\label{ch:literature}

\section{Introduction}
\label{sec:lit-introduction}

This chapter reviews the work that the proposed model builds on and compares with. It covers, in
order, the classical schedulers, the trainable models, the continuous-time memory models, the
Bayesian inference tools used in this work, and the cognitive-science studies that motivate some
modelling choices. The chapter ends with a comparison of the methods and a statement of the gap
that this work addresses.

Figure~\ref{fig:tk_lit_timeline} places the main sources on a single time line, from Ebbinghaus's
work in 1885 to the benchmark data used in this report. The figure is a chronology only. It does not
rank the methods.

\begin{figure}[htbp]
  \centering
  % tk_lit_timeline: literature milestones 1885--2026 (DIAGRAMS.md item 13)
% Label positions come from a greedy placement: no label overlaps another and
% no leader line runs through a label.
\begin{tikzpicture}[font=\footnotesize]
  \draw[arrow] (-0.2,0) -- (13.4,0);
  \draw[black!45, line width=0.4pt] (0.000,0) -- (0.000,+0.95);
  \node[event, minimum size=1.8mm] at (0.000,0) {};
  \node[anchor=south west, align=left, inner sep=1pt] at (0.000,+0.95) {Ebbinghaus \\ 1885};
  \draw[black!45, line width=0.4pt] (7.560,0) -- (7.560,+0.95);
  \node[event, minimum size=1.8mm] at (7.560,0) {};
  \node[anchor=south west, align=left, inner sep=1pt] at (7.560,+0.95) {Pimsleur \\ 1967};
  \draw[black!45, line width=0.4pt] (8.021,0) -- (8.021,-0.95);
  \node[event, minimum size=1.8mm] at (8.021,0) {};
  \node[anchor=north west, align=left, inner sep=1pt] at (8.021,-0.95) {Leitner \\ 1972};
  \draw[black!45, line width=0.4pt] (9.681,0) -- (9.681,+0.95);
  \node[event, minimum size=1.8mm] at (9.681,0) {};
  \node[anchor=south west, align=left, inner sep=1pt] at (9.681,+0.95) {SM-2 (Wozniak) \\ 1990};
  \draw[black!45, line width=0.4pt] (10.050,0) -- (10.050,-2.30);
  \node[event, minimum size=1.8mm] at (10.050,0) {};
  \node[anchor=north east, align=right, inner sep=1pt] at (10.050,-2.30) {Carter \& Kohn; \\ Fr\"uhwirth-Schnatter (FFBS), 1994};
  \draw[black!45, line width=0.4pt] (11.801,0) -- (11.801,-3.65);
  \node[event, minimum size=1.8mm] at (11.801,0) {};
  \node[anchor=north east, align=right, inner sep=1pt] at (11.801,-3.65) {Polson et al. \\ P\'olya-Gamma, 2013};
  \draw[black!45, line width=0.4pt] (12.078,0) -- (12.078,-5.00);
  \node[event, minimum size=1.8mm] at (12.078,0) {};
  \node[anchor=north east, align=right, inner sep=1pt] at (12.078,-5.00) {Settles \& Meeder \\ HLR, 2016};
  \draw[black!45, line width=0.4pt] (12.262,0) -- (12.262,+2.30);
  \node[event, minimum size=1.8mm] at (12.262,0) {};
  \node[anchor=south east, align=right, inner sep=1pt] at (12.262,+2.30) {Walsh et al. \\ PPE, 2018};
  \draw[black!45, line width=0.4pt] (12.631,0) -- (12.631,+3.65);
  \node[event, minimum size=1.8mm] at (12.631,0) {};
  \node[anchor=south east, align=right, inner sep=1pt] at (12.631,+3.65) {Ye et al. \\ FSRS, 2022};
  \draw[black!45, line width=0.4pt] (12.723,0) -- (12.723,+5.00);
  \node[event, minimum size=1.8mm] at (12.723,0) {};
  \node[anchor=south east, align=right, inner sep=1pt] at (12.723,+5.00) {Zhou et al. \\ HOUM, 2023};
  \draw[black!45, line width=0.4pt] (12.908,0) -- (12.908,+6.35);
  \node[event, minimum size=1.8mm] at (12.908,0) {};
  \node[anchor=south east, align=right, inner sep=1pt] at (12.908,+6.35) {FSRS-6; anki-revlogs-10k \\ 2024--26};
  \draw[black!45, line width=0.4pt] (13.000,0) -- (13.000,+7.70);
  \node[event, minimum size=1.8mm] at (13.000,0) {};
  \node[anchor=south east, align=right, inner sep=1pt] at (13.000,+7.70) {This work \\ 2026};
\end{tikzpicture}

  \caption{Milestones in the literature on memory and spaced-repetition models, from 1885 to 2026.
    The figure orders the sources in time. It does not rank them.}
  \label{fig:tk_lit_timeline}
\end{figure}

The review is organised by the role each family plays in this work. The classical schedulers are the
historical baseline. The trainable models are the baselines that the proposed model is compared with
directly. The continuous-time models are the closest in structure to the proposed model. The Bayesian
inference tools are the means by which the proposed model is fitted. The cognitive-science studies are the
evidence that motivates some of the modelling choices. Each section ends with the consequence for this work.

Two points of terminology are used throughout. A \emph{scheduler} is the rule that decides when an
item is shown again. A \emph{recall model} is a model that predicts the probability that a learner
recalls an item at a given time. The schedulers reviewed in Section~\ref{sec:lit-classical} are
rules. The models in Sections~\ref{sec:lit-trainable} to \ref{sec:lit-bayesian} are fitted to review
logs and produce recall probabilities.

\section{Classical Schedulers}
\label{sec:lit-classical}

The earliest spaced-repetition schedulers were fixed rules, designed by hand and applied to every
learner in the same way. The benefit of spacing was known well before these rules were written.
Ebbinghaus \cite{ebbinghaus1885memory} showed that retention of newly learned material falls with
time, and that repetition slows the fall. The schedulers below turn this observation into a rule for
when to review.

Pimsleur \cite{pimsleur1967} proposed a memory schedule in which the intervals between recalls of an
item grow in graded steps. The schedule is simple to follow, but it is one fixed sequence for all
items and learners.

Leitner \cite{leitner1972lernen} introduced a box system for flashcards (Figure~\ref{fig:tk_leitner}).
Cards are kept in numbered boxes. A correct answer moves the card to the next box, and a wrong answer
returns it to the first box. Each box is reviewed less often than the one before it, so well-known
cards are seen less frequently. The system needs no computer and is easy to explain, but it has no
parameters that depend on the learner, and it gives no probability that an answer is correct.

\begin{figure}[htbp]
  \centering
  % tk_leitner: the Leitner box system (DIAGRAMS.md item 14)
\begin{tikzpicture}[font=\small]
  \foreach \i/\d in {1/1,2/2,3/4,4/8,5/16} {
    \pgfmathsetmacro{\xc}{(\i-1)*2.8}
    \node[box, minimum width=2.1cm, minimum height=1.5cm, text width=1.9cm] (b\i) at (\xc,0)
      {Box \i\\[2pt]{\footnotesize review every\\\d\ day\ifnum\d>1 s\fi}};
  }
  % correct answers move a card to the next box (green)
  \foreach \i/\j in {1/2,2/3,3/4,4/5} {
    \draw[-{Stealth[length=2.2mm]}, line width=0.8pt, draw=ours!85!black] (b\i.east) -- (b\j.west);
  }
  \node[font=\footnotesize, text=ours!85!black, align=center] at (1.4,1.15) {correct $\rightarrow$ next box};
  % wrong answers send a card back to Box 1 (red), from every box
  \draw[line width=0.8pt, draw=hlr!85!black] (b2.south) -- (b2.south |- 0,-1.6);
  \foreach \i in {3,4,5} { \draw[line width=0.8pt, draw=hlr!85!black] (b\i.south) -- (b\i.south |- 0,-1.6); }
  \draw[-{Stealth[length=2.2mm]}, line width=0.8pt, draw=hlr!85!black] (b5.south |- 0,-1.6) -- (-0.0,-1.6) -- (b1.south);
  \node[font=\footnotesize, text=hlr!85!black, anchor=north] at (5.6,-1.6) {wrong $\rightarrow$ Box 1 (from every box)};
\end{tikzpicture}

  \caption{The Leitner box system. A correct answer moves a card to the next box, and a wrong answer
    returns it to Box 1 from any box. Boxes with higher numbers are reviewed less often. The intervals of
    1, 2, 4, 8 and 16 days are an illustrative doubling schedule, not a schedule taken from the source.}
  \label{fig:tk_leitner}
\end{figure}

Wozniak \cite{wozniak1990sm2} described the SM-2 algorithm in a master's thesis. SM-2 keeps a
per-item interval and an ease factor. After each review, the learner's grade changes both quantities
through fixed update rules. SM-2 is widely used as a reference point in the spaced-repetition
literature. It adapts to each item, but its update rules are fixed in advance rather than fitted to
data, and it does not model uncertainty about the learner's state.

The classical schedulers persist for practical reasons. Their rules are short enough to state in a few
lines, they can be checked by hand, and they require no training data. A rule that is transparent to the
learner has value of its own, because the learner can see why an item is due. These are the strengths that a
data-driven model must match before it can replace a rule.

Taken together, the classical schedulers share three limits. Their rules do not estimate parameters
for a particular learner. They do not produce calibrated probabilities of recall. They do not
represent the state of memory explicitly, so they cannot say how the state changes between reviews.

\section{Trainable Models}
\label{sec:lit-trainable}

\paragraph{Half-life regression.}
Settles and Meeder \cite{settles2016trainable} introduced half-life regression (HLR). The model
assumes that recall probability decays exponentially with elapsed time $\Delta t$, with a half-life
$h$ that depends on the learner's history through a weighted sum of features:
\begin{equation}
  p = 2^{-\Delta t/h}, \qquad h = 2^{\theta^\top x},
  \label{eq:hlr}
\end{equation}
where $\theta$ is a weight vector and $x$ contains features of the review history. The features are
transformed counts of correct and incorrect responses. Settles and Meeder minimise a squared loss on the predicted recall
probability, plus a weighted squared loss on the half-life, with L2 regularisation. They optimise it with
AdaGrad, a stochastic gradient method. The model was developed for review data from the Duolingo
language-learning platform.
Settles and Meeder report an AUC of 0.538 for HLR on that data. The HLR used in this work is a port of
the official code. It uses the count features and a bias term, and it has no lexeme features.

The main difference between HLR and the later models lies in how the time gap enters. In HLR, elapsed time
appears only in the exponent of Equation~\eqref{eq:hlr}, so the half-life sets the whole shape of the decay.
The learner's history affects the half-life through the count features, but it does not change the shape of
the curve. FSRS uses elapsed time together with a stability value that changes at each review. The present
model goes further, because its state moves between reviews and its mean depends on the previous outcome.

The HLR AUC reported by Settles and Meeder is not directly comparable with the AUC obtained in this work.
The two evaluations use different data, different splits and different event rules. On the benchmark used
here, the HLR port has an AUC of 0.555 (Chapter~\ref{ch:results}), which is a separate result on a separate
dataset.

HLR is a useful baseline because it is simple and trainable. Its decay curve is fixed for a given
history, so a learner's state cannot change in a way that the curve does not already describe. It gives
point predictions without uncertainty. In the benchmark reported in Chapter~\ref{ch:results}, its
mean bias changes sign between short and long gaps, which is the pattern that motivated this work
(Section~\ref{sec:intro-motivation}).

\paragraph{Predictive performance equation.}
Walsh et al.\ \cite{walsh2018ppe} evaluate computational models of the spacing effect. Their
evaluation includes the predictive performance equation (PPE), a model of learning and forgetting
over repeated practice. The paper assesses both the theoretical adequacy of such models and their
applied potential. It is a review of the modelling approach and does not supply data for the present
work. The PPE is reviewed for context only. It is not compared with the proposed model.

\paragraph{FSRS.}
Ye, Su and Cao \cite{ye2022stochastic} formulate spaced-repetition scheduling as a stochastic
shortest-path problem and propose the Free Spaced Repetition Scheduler (FSRS). FSRS couples a memory
model with a scheduling rule. Its memory model keeps two latent quantities for each card, a
stability and a difficulty, and predicts recall from the elapsed time and the stability. The parameters
of the model can be set to defaults or fitted to a learner's own review history. The fitting uses the
optimiser of the py-fsrs library on the training reviews. A learner with too few informative training
reviews falls back to the default parameters. The open-source
implementation \cite{pyfsrs} and the public benchmark dataset \cite{ankirevlogs10k} make FSRS
reproducible. The current version, FSRS-6, has 21 parameters. In this work, FSRS-6 is run in two
forms: with the default parameters, and with parameters fitted per user on the training split.

FSRS is the strongest of the established spaced-repetition models in this comparison (the later Experiment 4 adds a stronger tree model), and it is the reference against which the present
model is judged. Its memory state is updated deterministically once the parameters are fixed. It gives
point predictions rather than a posterior distribution.

\paragraph{Outcome coding.}
The trainable models differ in what they take as the outcome. FSRS uses a graded response with four levels,
which carries more information than a yes-or-no answer. HLR uses counts of correct and incorrect responses.
The proposed model uses a binary outcome: a review counts as recalled when its rating is greater than one. This
choice is a limitation. It discards the distinction between the grades above the threshold, and it is the
reason the proposed model cannot be compared with FSRS-6 on its graded predictions. The binary coding is kept
because the logistic observation model and the P\'olya--Gamma augmentation are defined for binary outcomes, and
because it allows a direct comparison of the recall probability with the observed recall indicator.

\paragraph{Static logistic baseline.}
Experiment 1 uses a static logistic model, LagLogit+$\Delta t$, as a second baseline. Its inputs are
the outcome of the previous review of the same card and the elapsed time. It has no latent state. It is
the simplest model in the comparison that uses the same information as the proposed model, and so it
isolates the effect of the latent state.

\paragraph{Feature-based and tree models.}
A different line of work does not model the memory state and predicts recall from features of the practice history.
DASH \cite{lindsey2014improving} and DAS3H \cite{choffin2019das3h} are logistic models of past practice across items and time windows.
In the Second Language Acquisition Modeling task on Duolingo data \cite{settles2018slam}, the winning entry was an ensemble
of LightGBM \cite{ke2017lightgbm} and a recurrent network \cite{osika2018slam}, and KARL \cite{shu2024karl} predicts recall with
knowledge-aware retrieval and learned representations. The open \texttt{srs-benchmark} \cite{srsbenchmark} compares spaced-repetition
models on about ten thousand Anki users. To our knowledge from that repository, its published models include no tree model, and its best published model, RWKV-Instant, uses
history across cards \cite{srsbenchmark}. A signal of the learner's overall ability across cards is therefore not new. Experiment 4 uses a gradient-boosted tree model \cite{friedman2001greedy} of this kind as
a strong, non-Bayesian reference, and Experiment 5 runs it pooled under the benchmark's protocol on 1,000 of its users (Section~\ref{sec:results-exp5-srsb}).

\paragraph{Sequence models and online learner ratings.}
The strongest published models on the srs-benchmark are neural sequence models. RWKV \cite{peng2023rwkv} is a recurrent
architecture with the training efficiency of a Transformer, and the benchmark's RWKV models read every review of a user,
across all cards, with their timing, grades and answer times. A much older and lighter way of tracking a learner is the Elo
rating system. Pel\'anek \cite{pelanek2016elo} reviews its use in adaptive educational systems: each learner has an ability
and each item a difficulty, the predicted probability of success is a logistic function of their difference, and both are
nudged after every answer by the prediction error. Experiment~6 uses Elo ratings to analyse learners by strength, and an
Elo-style online update to adjust the GBDT's prediction for each learner (Section~\ref{sec:method-gbdt-ext}).

\section{Continuous-Time Memory Models}
\label{sec:lit-continuous}

Discrete-time models treat each review as a step in a sequence. Continuous-time models treat memory as a
process that evolves between reviews. The mathematical framework for such models is the linear
stochastic differential equation (SDE). S\"arkk\"a and Solin \cite{sarkka2019applied} give a
comprehensive treatment of SDEs and their use in state estimation. The Ornstein--Uhlenbeck (OU) process
in Equation~\eqref{eq:ou} of Chapter~\ref{ch:introduction} is a linear SDE with mean reversion. For a
scalar state its transition over a gap of length $\Delta t$ has a closed form, so the process can be
evaluated exactly between reviews. For a state of dimension greater than one, the transition is
computed with the matrix exponential, a method described by Van Loan \cite{vanloan1978computing}. The
proposed model has a scalar state, so the closed form is sufficient and the matrix exponential is not
needed.

For the OU process the transition over a gap is an autoregression of order one,
\begin{align}
  \psi_t &= \phi_t\,\psi_{t-1} + (1-\phi_t)\,\mu_t + w_t, \qquad w_t \sim \mathcal{N}(0, q_t),
  \label{eq:ar1}\\
  \phi_t &= e^{-\lambda \Delta t_t}, \qquad
  q_t = \frac{\sigma^2\bigl(1 - e^{-2\lambda \Delta t_t}\bigr)}{2\lambda}.
\end{align}
where $\mu_t$ is the equilibrium mean at review $t$. The coefficient and the noise variance both depend on
the gap, so a model with irregular gaps must recompute them for each review. A discrete-time model with a
fixed step uses one coefficient for every gap, which cannot represent the difference between a gap of a
few hours and a gap of several months. Chapter~\ref{ch:methodology} derives Equation~\eqref{eq:ar1} in full.

Zhou, Tejero-Cantero and Wu \cite{zhou2023houm} proposed a hierarchical Ornstein--Uhlenbeck model (HOUM)
of learning and memory. The paper was presented at CCN in 2023. This report cites it from its title and venue only. The inference method used in HOUM is not
described here. The proposed model is related in its use of an OU latent state. Its observation model,
its lagged-outcome term and its inference scheme are specified in Chapter~\ref{ch:methodology}, and no
further comparison with HOUM is made.

The rate of mean reversion $\lambda$ has a direct interpretation. When $\lambda$ is close to zero, the state
hardly reverts to its mean over a typical gap, and the process behaves like a random walk. When $\lambda$ is
large, the state returns quickly to its equilibrium, so that a previous review has little effect on the next
one. The equilibrium mean sets the level to which the state returns. In the proposed model it depends on a
global intercept, a learner effect, a card effect and the outcome of the previous review. The rate itself may
depend on the number of earlier successful recalls through a term $\kappa$. These choices are specified in
Chapter~\ref{ch:methodology}. Their interaction contributes to the identifiability problems discussed in
Chapter~\ref{ch:results}.

Continuous-time models answer a question that the trainable models leave open. They make the state
change explicit between reviews, so that the effect of a long gap is a property of the dynamics rather
than of a fitted curve. The price is that the state must be inferred from data, and the inference is the
main difficulty addressed in the next section.

\section{Bayesian State-Space Inference}
\label{sec:lit-bayesian}

A state-space model consists of a state equation that describes how a latent state evolves and an
observation equation that links the state to the data. West and Harrison \cite{west1997bayesian} present
the dynamic linear model (DLM) as a framework for Bayesian forecasting with such models. For linear
Gaussian models, the Kalman filter gives the posterior distribution of the state exactly, given the
parameters. The observation model in this work is not linear Gaussian, because the outcome is binary.
Several tools are therefore combined.

\paragraph{P\'olya--Gamma augmentation.}
Polson, Scott and Windle \cite{polson2013pg} show that the logistic likelihood can be written as a
Gaussian likelihood conditional on latent P\'olya--Gamma variables. For an outcome $y \in \{0,1\}$ and
a logit $\psi$,
\begin{equation}
  \frac{e^{y\psi}}{1 + e^{\psi}}
  = 2^{-1} e^{(y - \tfrac{1}{2})\psi}
    \int_0^\infty e^{-\omega\psi^2/2}\, p(\omega \mid 1, 0)\,\mathrm{d}\omega,
  \qquad \omega \sim \mathrm{PG}(1, 0),
  \label{eq:pg}
\end{equation}
where $\mathrm{PG}(1, 0)$ denotes the P\'olya--Gamma distribution. Conditional on the auxiliary variable
$\omega$, the outcome acts like a Gaussian observation of the logit with pseudo-response
$z = (y - \tfrac{1}{2})/\omega$ and variance $1/\omega$. This makes the conditional updates of the state
and the regression coefficients conjugate. Polson et al.\ give an exact sampler for the augmenting
variables. The present work uses a truncated series instead, as Chapter~\ref{ch:methodology} explains.

\paragraph{Kalman recursions for a scalar state.}
For a scalar state with a Gaussian transition and a Gaussian observation, the Kalman filter has two steps at
each review. The prediction step propagates the current mean and variance through the transition. The update
step combines the prediction with the new observation, weighted by their variances. The pseudo-observation from
the P\'olya--Gamma representation fits into the update step, because it is Gaussian conditional on the augmenting
variable. The filter therefore needs no matrix inversion when the state is scalar, which is one reason the
scalar case is convenient.

\paragraph{Forward filtering, backward sampling.}
When the observation model is conditionally Gaussian, the state trajectory can be sampled jointly from
its posterior. Carter and Kohn \cite{carter1994gibbs} and Fr\"uhwirth-Schnatter \cite{fruhwirth1994data}
proposed forward-filtering backward-sampling (FFBS) for this purpose. A forward Kalman filter computes the
filtered distributions, and a backward pass draws the trajectory from them. For a scalar state the
recursions are short, and their cost grows linearly with the number of reviews. The Gibbs sampler then
alternates between the state trajectory and the other parameters.

\paragraph{Compilation and sampling.}
The FFBS recursions are sequential and contain many small operations, so the implementation runs them in
a compiled kernel. Lam, Pitrou and Seibert \cite{lam2015numba} describe Numba, a just-in-time compiler
for Python based on LLVM. Some parameters, notably the rate of mean reversion, do not have conjugate
conditional distributions. For these, a random-walk Metropolis--Hastings step is used inside the Gibbs
sweep. The combined procedure is Metropolis-within-Gibbs and is not an exact Gibbs sampler.

\paragraph{Convergence diagnostics.}
A sampler must be checked before its output is used. Gelman and Rubin \cite{gelman1992inference} introduced
the potential scale reduction factor, $\hat{R}$, which compares the variance within chains with the
variance between chains. Vehtari et al.\ \cite{vehtari2021rank} improve it with rank normalisation and
folding, and they recommend a threshold of 1.01 together with an effective sample size (ESS) that is large
enough for the quantities of interest. This work uses the same two checks, with an ESS of at least 400. A
sampler that fails them does not provide reliable posterior summaries, and the scores based on it are
treated as provisional.

\paragraph{Hierarchical priors.}
A hierarchical model shares information across groups through a common prior. Each learner's effect is drawn
from a distribution whose spread is itself estimated from the data. When the data for a learner are sparse, the
estimate is pulled toward the population mean, and when the data are plentiful, it follows the learner. This
partial pooling is the reason hierarchical models suit learners with different amounts of history. In the
proposed model the spreads of the learner and card effects receive inverse-gamma priors, which are conjugate
to the normal distribution of the effects. For a scalar model this is the natural choice, and it avoids the
matrix-valued priors that a multi-dimensional state would require.

\paragraph{Identifiability.}
Some parameters enter a model only through a sum. In a model with an intercept and additive learner and card
effects, a constant can be moved between the intercept and the effects without changing any likelihood value.
Such a model is not identified, and a sampler will move along the flat direction without converging. Remedies
include a sum-to-zero constraint, a reparameterisation that removes the redundant direction, and informative
priors. Each remedy changes what the parameters mean, so the choice should be stated. A related problem occurs
when a parameter is informed only through a combination of two others. This report returns to both problems in
Chapter~\ref{ch:results}, where the sampler diagnostics are examined.

\paragraph{Alternatives to sampling.}
Variational methods and other deterministic approximations to the posterior are faster than sampling, and they
are attractive for large datasets. Their approximation error is difficult to measure without a reference
posterior, and a sampler with convergence checks gives a way to see when the approximation is poor. For this
reason the present work uses a sampler, and it accepts the higher cost. Chapter~\ref{ch:methodology} discusses
the trade-off in more detail.

\paragraph{Interweaving.}
Yu and Meng \cite{yu2011interweaving} propose ancillarity--sufficiency interweaving (ASIS), which
alternates between two parameterisations of a model to improve mixing. An interweaving kernel for this
model was designed and implemented as a pilot. No results from it are available, and it is listed as
future work in Chapter~\ref{ch:conclusion}.

\section{Cognitive-Science Evidence}
\label{sec:lit-cognitive}

The studies in this section are cited to motivate modelling choices. They do not test the proposed model, and
none of them is used to fix a parameter or to validate a result. Each study is described from its title and
abstract.

\paragraph{Individual differences in retrieval practice.}
Fordyce et al.\ \cite{fordyce2025individual} study whether working-memory capacity changes the benefit of
retrieval practice. Across all three experiments, retrieval outperformed restudying, and the benefit held
regardless of individual differences in working-memory ability. The authors varied the learning procedure and
the scoring of working memory to test whether either would produce a moderation effect, and neither did.
De Lima and Buratto \cite{delima2024retrieval} review the same question. They find no consistent relationship
between individual-difference variables and the retrieval practice effect, and they attribute the difficulty of
interpreting the literature to methodological heterogeneity across studies.

Neither source establishes that learners differ in how much they benefit from retrieval practice. Fordyce et al.\
do not support such variation, and the review finds none that is consistent. The per-user effects $\alpha_u$ and
$\gamma_u$ in the proposed model are therefore a modelling choice. They are justified by the general heterogeneity
of learners' data, not by evidence of differential benefit. The proposed model estimates them from the data and
adds no individual-difference covariates.

\paragraph{Mental effort and load.}
Evans et al.\ \cite{evans2024cognitive} examine the relation between instructional strategies, cognitive
load and motivation in a survey of secondary-school students in Australia. Their outcomes are
load, motivation, engagement and achievement, not memory retention. Endres et al.\ \cite{endres2025developing}
develop a translingual scale that separates mental effort, which is invested deliberately, from mental load,
which is perceived passively. Neither study measures recall probability, and neither bears directly on the
model. They are relevant to a future version of the model that could use effort ratings, since the benchmark
data contain only recall outcomes and grades. Until such ratings are collected, the scales offer a
vocabulary for the design of that future study rather than an input to the proposed model.

\paragraph{Retrieval practice in applied settings.}
Serra et al.\ \cite{serra2025retrieval} review retrieval practice in the health professions, where it is used
for course and licensing examinations. The review supports the choice of recall as the outcome of interest in
applied learning. It does not provide a model or parameters.

\paragraph{Spacing in non-neural cells.}
Kukushkin et al.\ \cite{kukushkin2024massed} show a massed-spaced effect in immortalised, non-neural human cell
lines. Spaced pulses of stimulation produce a stronger and more sustained response than a single massed pulse,
and the effect depends on the ERK and CREB signalling pathways. The authors suggest that canonical memory features can be embedded in signalling dynamics without neural
circuits. The study concerns cell lines, not learners, and it does not transfer any parameter to human
recall. It is cited only to note that spacing effects have been reported outside neural circuits. It is at
most a loose analogy for describing spacing through the time between events, as the proposed model does, and
it does not support any feature of the proposed model.

\paragraph{Consolidation across sleep.}
Guttesen et al.\ \cite{guttesen2026memory} review the resource reallocation hypothesis, which proposes that
sleep consolidation frees hippocampal capacity for next-day learning. They conclude that the evidence for this
hypothesis is not clear. A night between two reviews is therefore a mixture of sleep-related consolidation and
forgetting. The proposed model treats all gaps in the same way, which is a limitation discussed in
Chapter~\ref{ch:conclusion}.

\paragraph{Reinforcement learning and memory.}
Lee and Jung \cite{lee2025memory} present a theoretical perspective in which the hippocampal CA3 region generates
activity patterns and the CA1 region evaluates and reinforces them, a scheme that parallels the Dyna
reinforcement-learning algorithm. The perspective contains no data and provides no parameters for the present
model. It is cited only as an example of a theory that separates generating and evaluating components of memory.
It does not support any particular feature of the proposed model.

\paragraph{Further sources.}
Three further sources are cited from their titles and venues only, and no details of their methods or results are
used. They are Zou et al.\ \cite{zou2025cell}, \emph{Cell Reports}; Yang et al.\ \cite{yang2025commbio},
\emph{Communications Biology}; and Guttesen et al.\ \cite{guttesen2025cognition}, \emph{Cognition}.

Taken together, these studies are consistent with three design choices: a per-learner component, a recall
outcome as the target, and a time-dependent state. None of them establishes that any of these choices is correct. Their
value lies in showing that the choices are reasonable given what is known about learners, and in identifying
the data that a future study would need to collect, such as effort ratings.

\section{Comparative Summary}
\label{sec:lit-summary}

Accuracy measures differ across the studies reviewed in this chapter, and a direct comparison of their
published numbers is therefore difficult. The benchmark in this report therefore uses several measures on the
same events: log loss, Brier score, AUC, expected calibration error, and the mean bias within each gap bin.
Chapter~\ref{ch:results} defines them.

Table~\ref{tab:lit_summary} summarises the methods in the trainable and continuous-time families. Figure~\ref{fig:tk_taxonomy}
places them in a taxonomy. The table records each method's strength and limitation as they bear on the present
problem. The entry for HOUM is limited to what its title and venue support.

\begin{table}[htbp]
  \centering
  \footnotesize\setlength{\tabcolsep}{4pt}
  \caption{Summary of the main methods, their data, their strengths and limitations, and their relation to this work.}
  \label{tab:lit_summary}
  \footnotesize
  \begin{tabular}{@{}>{\raggedright\arraybackslash}p{2.5cm}>{\raggedright\arraybackslash}p{3.5cm}>{\raggedright\arraybackslash}p{3.8cm}>{\raggedright\arraybackslash}p{3.0cm}@{}}
    \toprule
    Source & Approach and data & Strength and limitation & Relation to this work \\
    \midrule
    Pimsleur (1967) & Graded interval rule; no data fitted & Simple; one fixed schedule for all learners & Background only \\
    Leitner (1972) & Box rule for flashcards; no data fitted & Simple and transparent; no learner parameters or probabilities & Background; Figure~\ref{fig:tk_leitner} \\
    SM-2, Wozniak (1990) & Per-item interval rule with ease factor & Adapts to each item; update rules are fixed, not fitted; no uncertainty & Reference scheduler \\
    HLR, Settles \& Meeder (2016) & Exponential decay with count-based half-life; learner review data from Duolingo & Trainable and simple; fixed curve, point estimates, AUC 0.538 on its own data & Baseline; ported and evaluated here \\
    PPE, Walsh et al.\ (2018) & Review of spacing models, including the PPE & Evaluates theoretical and applied value; no new data & Context \\
    FSRS, Ye et al.\ (2022) & Stochastic shortest-path scheduling with a stability and difficulty memory model; learner review logs & Strongest established model here (Exp.\ 2); point estimates, deterministic state & Baseline in two forms (default, optimised) \\
    HOUM, Zhou et al.\ (2023) & Hierarchical OU model of learning and memory & Continuous-time latent state; inference method and results not verified & Related; not compared in detail \\
    FFBS, Carter \& Kohn (1994); Fr\"uhwirth-Schnatter (1994) & Forward filtering and backward sampling for state-space models & Exact for conditionally Gaussian models; needs conjugate structure & Used for the latent state \\
    P\'olya--Gamma, Polson et al.\ (2013) & Auxiliary-variable representation of the logistic likelihood & Makes conditional updates conjugate; approximation needed in practice & Used; truncated approximation \\
    \bottomrule
  \end{tabular}
\end{table}

\begin{figure}[htbp]
  \centering
  % tk_taxonomy: families of memory and recall models (DIAGRAMS.md item 15), forest
\begin{forest}
  for tree={font=\footnotesize, draw=black!70, fill=black!4, rounded corners=3pt,
    inner sep=3pt, align=left, text width=4.7cm, grow=east,
    edge={draw=black!75, -{Stealth[length=2mm]}}, parent anchor=east, child anchor=west,
    s sep=4pt, l sep=12pt}
  [{\textbf{Memory / recall models}}, text width=4.0cm, fill=black!10
    [{\textbf{Heuristic schedulers}}, text width=4.0cm
      [Pimsleur (1967)]
      [Leitner (1972)]
      [SM-2 (1990)]
    ]
    [{\textbf{Trainable regression}}, text width=4.0cm
      [HLR]
      [PPE (Walsh et al.\ 2018)]
      [LagLogit + $\Delta t$]
    ]
    [{\textbf{State-based models}}, text width=4.0cm
      [FSRS-6 (deterministic 2-D state)]
      [HOUM (continuous-time OU)]
      [{Ours (Bayesian OU state-space)}, fill=ours!15, draw=ours!85!black]
    ]
    [{\textbf{Neural sequence models}}, text width=4.0cm
      [{RNN / LSTM, Transformers\\(mentioned only)}, text width=4.7cm]
    ]
  ]
\end{forest}

  \caption{Families of memory and recall models discussed in this chapter. The shaded box marks the present
    model. Neural sequence models are mentioned for completeness and are not reviewed.}
  \label{fig:tk_taxonomy}
\end{figure}

In short, the table shows three options. Interpretable rules are transparent but cannot be fitted to a
learner. Fitted point models, such as HLR and FSRS, predict well in practice, and FSRS is the stronger of the two
on the benchmark used here. A latent-state model gives a principled account of the state and its uncertainty,
but its inference is costly and must pass convergence checks before its predictions can be trusted. The present
work takes the third option and tests whether its extra structure pays off.

Table~\ref{tab:feature_comparison} compares the models on the features that matter most for this problem. The
entries follow the descriptions in the preceding sections. The table is qualitative and does not rank the
models. It does not include the Experiment~4 models (v2, v3 and the gradient-boosted tree model, Section~\ref{sec:results-exp4}). A model that scores well on every feature need not predict best, and the benchmark results in
Chapter~\ref{ch:results} are the test of prediction. The table shows which design choices each model makes.
Two rows deserve comment. The FSRS-6 row marks per-user parameters as present only when they are optimised,
because the default parameters are shared. The proposed model's row marks the uncertainty as present but
unverified in convergence, because a posterior that has not converged cannot be relied on for its spread. A cell marked as not verified is one whose
details were not checked against the source.

\begin{table}[htbp]
  \centering
  \footnotesize\setlength{\tabcolsep}{4pt}
  \caption{Feature comparison of SM-2, HLR, FSRS-6, HOUM and the proposed model.}
  \label{tab:feature_comparison}
  \footnotesize
  \begin{tabular}{@{}>{\raggedright\arraybackslash}p{1.9cm}>{\raggedright\arraybackslash}p{2.15cm}>{\raggedright\arraybackslash}p{2.15cm}>{\raggedright\arraybackslash}p{2.15cm}>{\raggedright\arraybackslash}p{2.15cm}>{\raggedright\arraybackslash}p{2.15cm}@{}}
    \toprule
    Model & Latent state & Continuous time & Uncertainty & Per-user & Graded ratings \\
    \midrule
    SM-2 & No (per-item interval) & No (day intervals) & No & No & Yes \\
    HLR & No & Yes (exponential in $\Delta t$) & No & No (global weights) & No (correct and wrong counts) \\
    FSRS-6 & Yes (stability, difficulty) & Yes & No (point estimates) & Yes, when optimised & Yes (four grades) \\
    HOUM & Yes (OU state) & Yes (OU) & Not verified & Not verified & Not verified \\
    Present model & Yes (scalar $\psi$) & Yes (OU) & Yes (posterior; sampler not converged) & Yes for $\alpha_u$, $\gamma_u$; $\lambda$ and $\kappa$ pooled & No (binary outcome) \\
    \bottomrule
  \end{tabular}
\end{table}

\section{Research Gap}
\label{sec:lit-gap}

The sources reviewed in this chapter suggest three gaps.

\begin{enumerate}
  \item \textbf{Point estimates.} The classical and trainable models give point predictions and do not
    represent the state of memory explicitly. A learner's state is either absent or fixed between reviews.
  \item \textbf{Untested continuous-time models.} The continuous-time memory models reviewed here are not
    tested on a public, leakage-free benchmark against the strongest established trainable model, to the knowledge of the
    sources examined in this chapter.
  \item \textbf{Uncertainty for the strongest established model.} FSRS-6 is fitted per user as a point estimate, and no
    posterior uncertainty is reported for it. The effect of this on the reliability of its predictions is not
    established by the sources reviewed.
\end{enumerate}

Evidence that would close these gaps is specific. For the first, a model must show that an explicit state
improves prediction, not only that it fits. For the second, the model must be scored on a public benchmark
with a leakage-free split and compared with the strongest existing model. For the third, the uncertainty of
the strongest model must be reported and checked against outcomes. A result on one of these points does not
settle the others.

This work addresses these gaps in part. It builds a continuous-time model with a scalar latent state and a
Bayesian inference scheme for binary outcomes. It evaluates the model on a public benchmark against FSRS-6 and
HLR, under the leakage-free protocol described in Chapter~\ref{ch:dataset}. In Experiment 2 on Anki the model is worse than
FSRS-6 in every gap bin; in Experiment 3 on Duolingo it has the lowest log loss of the five models of that experiment, with the caveats in Section~\ref{sec:results-exp3}. The model is clearly better than HLR on Anki.
The sampler did not converge, so the scores are provisional. A fourth question arose from these results: how much of the
error is due to the memory model at all. Experiment 4 therefore adds two MAP variants of the hierarchy and a gradient-boosted tree
model, and finds that the tree model, which has no memory equation, is the best predictor of the models compared on both datasets (Section~\ref{sec:results-exp4}). Chapter~\ref{ch:dataset} defines the comparison.
